% ══════════════════════════════════════════════════════════════
%  Chapter 12: Network Layer --- P2P Message Routing and Sync
% ══════════════════════════════════════════════════════════════
\chapter{Network Layer}
\label{ch:network}

\begin{learnbox}
\begin{itemize}
  \item Process peer-to-peer messages for blockchain synchronization
  \item Manage peer connections and handle discovery of new nodes
  \item Coordinate blockchain synchronization across a distributed network
  \item Understand the network protocol that enables decentralized consensus
\end{itemize}
\end{learnbox}

\lettrine{A}{} \index{blockchain}blockchain that cannot communicate with peers is just a local database. The \index{network layer}network layer is what turns isolated nodes into a \index{consensus}consensus system: it propagates new \index{transaction}transactions so \index{mining}miners can include them in \index{block}blocks, relays newly mined \index{block}blocks so all nodes converge on the same chain, and synchronizes new nodes that need to catch up with the canonical history. Every guarantee Bitcoin makes --- immutability, \index{double-spend}double-spend prevention, censorship resistance --- depends on this protocol layer working correctly.

\begin{notebox}[Scope]
This chapter implements a simplified P2P protocol for blockchain synchronization. We do not cover NAT traversal, relay nodes, the full Bitcoin protocol message set, or Tor/I2P privacy layers.
\end{notebox}

% ──────────────────────────────────────────────────────────────
\section{Chapter Map: What ``Network'' Means in This Repo}
\label{sec:ch12-chapter-map}

This repository's \index{peer-to-peer}P2P layer is deliberately small:

\begin{itemize}
  \item \textbf{\index{TCP}Transport}: \index{TCP}TCP streams
  \item \textbf{\index{encoding}Encoding}: \index{JSON}JSON via \code{\index{Rust!serialization}\index{serde}serde\_json}
  \item \textbf{\index{message}Message model}: \code{\index{Package}Package} enum + \code{\index{OpType}OpType} enum
  \item \textbf{\index{dispatcher}Dispatcher}: \code{\index{process\_stream}process\_stream(...)} (routes each \index{inbound}inbound \index{package}package to the right action)
  \item \textbf{\index{send}Send primitives}: \code{\index{send\_data}send\_data(...)} and typed wrappers (\code{\index{send\_inv}send\_inv}, \code{\index{send\_get\_data}send\_get\_data}, \code{\index{send\_block}send\_block}, \code{\index{send\_tx}send\_tx}, \code{\index{send\_version}send\_version}, \ldots)
  \item \textbf{\index{peer!bootstrap}Peer bootstrap}: \index{central-node}central-node concept + ``known \index{node}nodes'' exchange
\end{itemize}

The core implementation lives in:

\begin{itemize}
  \item \code{bitcoin/src/net/net\_processing.rs}
  \item \code{bitcoin/src/node/server.rs} (TCP accept loop + bootstrap wiring)
\end{itemize}

\begin{warnbox}[Implementation Note]
This implementation uses a simplified P2P protocol suitable for learning. It does not implement NAT traversal, relay nodes, or the full Bitcoin protocol message set. Production networking would require significantly more robust connection handling.
\end{warnbox}

% ──────────────────────────────────────────────────────────────
\section{The Minimal Protocol Loop}
\label{sec:ch12-protocol-loop}

Bitcoin uses a \textbf{\index{gossip protocol}gossip and fetch} strategy to minimize bandwidth: peers announce what they have by sharing only the hash (the \index{gossip protocol}gossip phase), peers request the full object by hash (the fetch phase), and peers deliver the full object.

\begin{minted}[fontsize=\footnotesize,linenos=false]{rust}
Peer A has object (tx or block)
  |
  | 1) announce (hash only)
  v
INV(op_type, [id]) -----> Peer B
                             |
                             | 2) request bytes
                             v
                          GETDATA -----> Peer A
                                           |
                                           | 3) send full bytes
                                           v
                                  (TX | BLOCK) -----> Peer B
                                                        |
                                                        | 4) hand to node
                                                        v
                                                 mempool / add_block
\end{minted}

This three-message pattern---announce, request, deliver---is the heartbeat of blockchain peer-to-peer communication.

% ──────────────────────────────────────────────────────────────
\section{Message Model: The Wire-Level Vocabulary}
\label{sec:ch12-message-model}

Before reading any runtime flow, we need to understand what messages exist and what fields they carry.

\begin{listing}[htbp]
\caption{Network globals and message type enums}
\label{lst:ch12-message-model}
\begin{minted}[fontsize=\footnotesize]{rust}
pub const NODE_VERSION: usize = 1;

// Discriminator: tells receiver what type
// (Tx or Block) the hash refers to.
#[derive(Debug, Serialize, Deserialize)]
pub enum OpType { Tx, Block }

#[derive(Debug, Serialize, Deserialize)]
pub enum MessageType {
    Error, Success, Info, Warning, Ack
}

// Admin operations (not part of P2P protocol)
#[derive(Debug, Serialize, Deserialize)]
pub enum AdminNodeQueryType {
    GetBalance { wlt_address: String },
    GetAllTransactions,
    GetBlockHeight,
    MineEmptyBlock,
    ReindexUtxo,
}

// Wire envelope: every JSON message is one Package
#[derive(Debug, Serialize, Deserialize)]
pub enum Package {
    Block { addr_from: SocketAddr, block: Vec<u8> },
    GetBlocks { addr_from: SocketAddr },
    GetData {
        addr_from: SocketAddr,
        op_type: OpType,
        id: Vec<u8>
    },
    Inv {
        addr_from: SocketAddr,
        op_type: OpType,
        items: Vec<Vec<u8>>
    },
    Tx { addr_from: SocketAddr, transaction: Vec<u8> },
    SendBitCoin {
        addr_from: SocketAddr,
        wlt_frm_addr: String,
        wlt_to_addr: String,
        amount: i32,
    },
    KnownNodes { addr_from: SocketAddr, nodes: Vec<SocketAddr> },
    Version {
        addr_from: SocketAddr,
        version: usize,
        best_height: usize
    },
    Message {
        addr_from: SocketAddr,
        message_type: MessageType,
        message: String,
    },
    AdminNodeQuery {
        addr_from: SocketAddr,
        query_type: AdminNodeQueryType
    },
}
\end{minted}
\end{listing}

% ──────────────────────────────────────────────────────────────
\section{TCP Server Loop: How Nodes Accept Connections}
\label{sec:ch12-server-loop}

This is the runtime entry point for inbound peer connections. It binds a TCP listener, does initial bootstrap, and then spawns one task per accepted connection.

\begin{listing}[htbp]
\caption{TCP server setup and bootstrap}
\label{lst:ch12-server-setup}
\begin{minted}[fontsize=\footnotesize]{rust}
#[derive(Debug, Clone)]
pub struct Server {
    node_context: NodeContext,
}

impl Server {
    pub fn new(blockchain: NodeContext) -> Server {
        Server { node_context: blockchain }
    }

    pub async fn run_with_shutdown(
        &self,
        addrs: &SocketAddr,
        connect_nodes: HashSet<ConnectNode>,
        mut shutdown:
            tokio::sync::broadcast::Receiver<()>,
    ) {
        let listener = TcpListener::bind(addrs)
            .await
            .expect("TcpListener bind error");
        info!("Server listening on {:?}",
            listener.local_addr());

        if !addrs.eq(&CENTERAL_NODE) {
            let best_height = self
                .node_context
                .get_blockchain_height()
                .await
                .expect("Blockchain read error");
            send_version(&CENTERAL_NODE, best_height)
                .await;
        } else {
            let remote_nodes: HashSet<SocketAddr> =
                connect_nodes
                    .iter()
                    .filter(|node| node.is_remote())
                    .map(|node| node.get_addr())
                    .collect();

            GLOBAL_NODES
                .add_nodes(remote_nodes.clone())
                .expect("Global nodes add error");
        }

        // Serve incoming connections with
        // graceful shutdown.
        loop {
            tokio::select! {
                _ = shutdown.recv() => {
                    info!(
                        "Network server shutdown signal received"
                    );
                    break;
                }
                accept_res = listener.accept() => {
                    match accept_res {
                        Ok((stream, _peer)) => {
                            let blockchain =
                                self.node_context.clone();
                            tokio::spawn(async move {
                                match stream.into_std() {
                                    Ok(std_stream) => {
                                        let _ = std_stream
                                            .set_nonblocking(false);
                                        if let Err(e) =
                                            net_processing
                                            ::process_stream(
                                                blockchain,
                                                std_stream,
                                            ).await {
                                            error!(
                                                "Serve error: {}",
                                                e
                                            );
                                        }
                                    }
                                    Err(e) => {
                                        error!(
                                            "Stream conversion error: {}",
                                            e
                                        );
                                    }
                                }
                            });
                        }
                        Err(e) => {
                            error!("Accept error: {}", e);
                        }
                    }
                }
            }
        }
    }
}
\end{minted}
\end{listing}

% ──────────────────────────────────────────────────────────────
\section{Message Dispatcher: Routing Inbound Messages}
\label{sec:ch12-dispatcher}

The dispatcher is the heart of message routing. It reads JSON from the TCP stream, deserializes into a \code{Package}, and routes based on the variant.

\begin{listing}[htbp]
\caption{Message dispatcher (simplified)}
\label{lst:ch12-dispatcher}
\begin{minted}[fontsize=\footnotesize]{rust}
pub async fn process_stream(
    node_context: NodeContext,
    stream: std::net::TcpStream,
) -> Result<()> {
    let reader = std::io::BufReader::new(
        stream.try_clone()?);

    // JSON message stream
    let mut deserializer =
        serde_json::Deserializer::from_reader(reader);

    // For each JSON message from the stream...
    loop {
        match serde_json::Value::deserialize(
            &mut deserializer) {
            Ok(value) => {
                // Parse the wire envelope
                match serde_json::from_value::<Package>(
                    value) {
                    Ok(pkg) => {
                        // Dispatch by variant
                        if let Err(e) =
                            match pkg {
                            Package::Tx { addr_from, transaction } => {
                                let tx = Transaction
                                    ::deserialize(&transaction)?;
                                node_context
                                    .process_transaction(
                                        &addr_from,
                                        tx
                                    ).await
                                    .map(|_| ())
                            }
                            Package::Block {
                                addr_from, block
                            } => {
                                let blk = Block
                                    ::deserialize(&block)?;
                                node_context
                                    .add_block(&blk)
                                    .await
                            }
                            Package::Inv {
                                addr_from,
                                op_type,
                                items
                            } => {
                                handle_inv(
                                    &node_context,
                                    &addr_from,
                                    op_type,
                                    items
                                ).await
                            }
                            Package::GetData {
                                addr_from,
                                op_type,
                                id
                            } => {
                                handle_get_data(
                                    &node_context,
                                    &addr_from,
                                    op_type,
                                    id
                                ).await
                            }
                            Package::Version {
                                addr_from,
                                version,
                                best_height
                            } => {
                                // Record peer info
                                GLOBAL_NODES
                                    .add_node(addr_from)
                                    .map(|_| ())
                            }
                            _ => Ok(()),
                        } {
                            error!("Handler error: {}", e);
                        }
                    }
                    Err(e) => {
                        error!("Deserialization error: {}", e);
                        break;
                    }
                }
            }
            Err(e) if e.is_eof() => {
                // Peer disconnected
                break;
            }
            Err(e) => {
                error!("JSON read error: {}", e);
                break;
            }
        }
    }
    Ok(())
}
\end{minted}
\end{listing}

% ──────────────────────────────────────────────────────────────
\section{Send Primitives: Broadcasting to Peers}
\label{sec:ch12-send-primitives}

The send primitives are typed wrappers around a base \code{send\_data(...)} function that opens a fresh TCP connection and sends one JSON message.

\begin{listing}[htbp]
\caption{Send primitives for inventory and block relay}
\label{lst:ch12-send}
\begin{minted}[fontsize=\footnotesize]{rust}
pub async fn send_inv(
    addr: &SocketAddr,
    op_type: OpType,
    items: &[Vec<u8>],
) {
    let pkg = Package::Inv {
        addr_from: GLOBAL_CONFIG.get_node_addr(),
        op_type,
        items: items.to_vec(),
    };
    let _ = send_data(addr, pkg).await;
}

pub async fn send_get_data(
    addr: &SocketAddr,
    op_type: OpType,
    id: Vec<u8>,
) {
    let pkg = Package::GetData {
        addr_from: GLOBAL_CONFIG.get_node_addr(),
        op_type,
        id,
    };
    let _ = send_data(addr, pkg).await;
}

pub async fn send_block(
    addr: &SocketAddr,
    block: &Block,
) -> Result<()> {
    let pkg = Package::Block {
        addr_from: GLOBAL_CONFIG.get_node_addr(),
        block: block.serialize()?,
    };
    send_data(addr, pkg).await
}

pub async fn send_tx(
    addr: &SocketAddr,
    tx: &Transaction,
) -> Result<()> {
    let pkg = Package::Tx {
        addr_from: GLOBAL_CONFIG.get_node_addr(),
        transaction: tx.serialize()?,
    };
    send_data(addr, pkg).await
}

pub async fn send_version(
    addr: &SocketAddr,
    best_height: usize,
) {
    let pkg = Package::Version {
        addr_from: GLOBAL_CONFIG.get_node_addr(),
        version: NODE_VERSION,
        best_height,
    };
    let _ = send_data(addr, pkg).await;
}

// Base send primitive: open connection, send message
async fn send_data(
    addr: &SocketAddr,
    pkg: Package,
) -> Result<()> {
    let mut stream = std::net::TcpStream::connect(addr)?;
    let _ = serde_json::to_writer(&mut stream, &pkg);
    let _ = stream.flush();
    Ok(())
}
\end{minted}
\end{listing}

% ──────────────────────────────────────────────────────────────
\section{Peer Discovery: Bootstrap and Known Nodes}
\label{sec:ch12-peer-discovery}

The network layer maintains a peer set and exchanges known nodes to bootstrap new peers. The central node concept provides an initial meeting point.

\begin{notebox}[Bootstrap Strategy]
Nodes start with a hardcoded central node address. They announce their presence by sending a \code{Version} message. The central node collects these addresses and redistributes them via \code{KnownNodes}. This provides the initial peer discovery mechanism without requiring a separate DHT or seed server.
\end{notebox}

% ──────────────────────────────────────────────────────────────
\section{Trade-offs: Transport Layer Architecture}
\label{sec:ch12-transport-tradeoffs}

The network layer currently uses \code{std::net::TcpStream} for compatibility with existing JSON deserialization code. This choice trades scaling for simplicity:

\begin{itemize}
  \item \textbf{Advantage}: Direct compatibility with \code{serde\_json::Deserializer::from\_reader(...)}.
  \item \textbf{Disadvantage}: Potential to block Tokio worker threads under high peer load.
\end{itemize}

A future migration to \code{tokio::net::TcpStream} would require:

\begin{enumerate}
  \item Moving from \code{serde\_json::from\_reader(...)} to async frame parsing
  \item Implementing message framing (length-delimited or newline-delimited)
  \item Updating connection handling to use async I/O primitives
\end{enumerate}

This would enable better scalability for production networks with many peers.

% ──────────────────────────────────────────────────────────────
\section{Exercises}
\label{sec:ch12-exercises}

\begin{enumerate}
  \item \textbf{Message Flow Trace} --- Run two nodes locally and submit a transaction to one node. Trace the P2P messages exchanged as the transaction propagates to the second node. Document the message sequence: which message types are sent, in what order, and what each contains.

  \item \textbf{Partition Recovery Scenario} --- Imagine a network partition splits your nodes into two groups. Each group mines blocks independently for 5 minutes. When the partition heals, describe the chain synchronization process. Which chain wins, and what happens to transactions in the losing chain?
\end{enumerate}

% ──────────────────────────────────────────────────────────────
\section{Common Errors}
\label{sec:ch12-common-errors}

\begin{notebox}[Troubleshooting: Network Layer]
\textbf{``Connection refused'' when connecting to a peer} --- The peer node is not running, or it is listening on a different port/interface than expected. Verify the peer's bind address (check for \code{0.0.0.0} vs \code{127.0.0.1}) and that no firewall rules are blocking the port.

\textbf{Timeout during block synchronization} --- If a new node is catching up with a long chain, the default read timeout may be too short. Increase the TCP timeout or implement chunked block transfer. Also check that the sending peer is not blocked on a lock while the receiver waits.

\textbf{``Address already in use'' on startup} --- The previous node process did not release the TCP port. Wait a few seconds for the OS to reclaim it, or use \code{SO\_REUSEADDR} in the socket options.

\textbf{Peers connect but no blocks propagate} --- Check that the message dispatcher is routing \code{NewBlock} messages to the chainstate handler. A common mistake is registering the handler for \code{NewTransaction} but forgetting \code{NewBlock}.
\end{notebox}

% ──────────────────────────────────────────────────────────────
\begin{chaptersummary}
\item We built the peer-to-peer networking layer that enables blockchain nodes to communicate and synchronize.
\item We implemented message processing for block and transaction propagation across the network.
\item We designed peer connection management for discovery, handshake, and ongoing communication.
\item We established the network synchronization protocol that keeps all nodes' chain state consistent.
\item We understood the trade-offs between simplicity (blocking TCP) and scalability (async I/O).
\end{chaptersummary}

In the next chapter, we bring all these subsystems together under a unified \code{NodeContext} API that coordinates blockchain state, mempool, network, and mining.
